\section*{Coda: a graph can be fixed in parameters and dynamic in response}
\addcontentsline{toc}{section}{Coda: a graph can be fixed in parameters and dynamic in response}

The static--dynamic decomposition changes the comparison between family models and protein
Transformers.  Their difference is no longer summarized by saying that one is generative and the
other contextual.  Both can be interrogated through the same categorical response operator.  The
Potts model occupies the background-invariant limit; the nonlinear model supplies a mean operator
and a measurable field of deviations.

This common coordinate does not erase the distinction between the models.  A mean pLM response is
an estimator built from interventions on a trained function, not a parameter fitted from an MSA.
A small dynamic residual supports a static approximation under the audited background law, not a
universal physical coupling.  A large residual is evidence against that approximation and a clue
to higher-order organization, but its decomposition depends on the reference geometry.

The next question is what kind of complexity the residual represents.  Graph support, categorical
rank, and interaction order are often compressed into one informal notion of complexity.  They are
independent.  A single edge may carry a full-rank operator; a dense graph may use only a shared
low-dimensional family of maps; and nonlinear composition can generate higher-order terms even
when every local message was separately pair-centered.  The next chapter separates these axes.
